\scanpage{54}
\begin{center}4 - 37 -\hfill n\textsuperscript{o} 392\end{center}

lui aussi injectif. Si $I=\operatorname{Ker}(f\otimes1)$, on a $I_{\underline{p}}=\operatorname{Ker}((f\otimes1)_{\underline{p}})$, (chap. II, § 2, n\textsuperscript{o} 4, th. 1) ; on a donc $I_{\underline{p}}=0$, d'où \uline{a fortiori} $I_{\underline{P}}=(I_{\underline{p}})_{\underline{P}}=0$ pour $\underline{P}\mid\underline{p}$, ce qui prouve le lemme, (n\textsuperscript{o} 4, prop. 9).

Revenons à la démonstration de (ii). Pour tout $A$-module de torsion de type fini $M$, posons $\varphi(M)=\unsure{\chi|\gamma}_B(M\otimes_A B)$ ; il résulte de (i) que si $M$ est pseudo-nul, on a $\varphi(M)=0$. D'autre part, considérons une suite exacte de $A$-modules de torsion de type fini :
\diagram{\begin{tikzcd}
0\arrow[r]&M_1\arrow[r]&M_2\arrow[r]&M_3\arrow[r]&0.
\end{tikzcd}}
Il résulte du lemme 4 que l'on a une suite exacte de $B$-modules :
\diagram{\begin{tikzcd}
0\arrow[r]&I\arrow[r]&M_1\otimes_A B\arrow[r]&M_2\otimes_A B\arrow[r]&M_3\otimes_A B\arrow[r]&0
\end{tikzcd}}
où $I$ est pseudo-nul. En utilisant le n\textsuperscript{o} 5, cor. de la prop. 10, on a donc $\varphi(M_2)=\varphi(M_1)+\varphi(M_3)$. On conclut donc de la prop. 11 du n\textsuperscript{o} 5 qu'il existe un homomorphisme $\unsure{\vartheta|\theta}:D(A)\longrightarrow D(B)$ tel que $\varphi(M)=\unsure{\vartheta|\theta}(\unsure{\chi|\gamma}_A(M))$ pour tout $A$-module de torsion de type fini $M$. Pour prouver que $\unsure{\vartheta|\theta}=i$, il suffit de montrer que $\varphi(A/\underline{p})=i(\underline{p})$ pour tout $\underline{p}\in P(A)$ ; or on a $(A/\underline{p})\otimes_A B=B/\underline{p}B$, et pour tout $\underline{P}\in P(B)$, $(B/\underline{p}B)_{\underline{P}}=B_{\underline{P}}/\underline{p}B_{\underline{P}}$ ; ce dernier module est $0$ si $\underline{P}$ n'est pas au-dessus de $\underline{p}$ ; si au contraire $\underline{P}\mid\underline{p}$, $B_{\underline{P}}/\underline{p}B_{\underline{P}}$ est un $B_{\underline{P}}$-module de longueur $e(\underline{P}/\underline{p})$ par définition de l'indice de ramification, (chap. VI, § 8, n\textsuperscript{o} 1) ; on a donc $\unsure{\chi|\gamma}_B(B/\underline{p}B)=\sum_{\underline{P}\mid\underline{p}}e_{\underline{P}/\underline{p}}.\underline{P}=i(\underline{p})$, ce qui achève de prouver (ii).

La formule (24) est immédiate, car $(M\otimes_A B)\otimes_B L=M\otimes_A L=(M\otimes_A K)\otimes_K L$ et la dimension sur $L$ de $(M\otimes_A K)\otimes_K L$ est égale à la dimension sur $K$ de $M\otimes_A K$. Pour démontrer (23), considérons un sous-module libre $H$ de $M$ tel que $Q=M/H$ soit un $A$-module de torsion. Appliquant comme ci-dessus le lemme 4, on a une suite exacte de $B$-modules :
\diagram{\begin{tikzcd}
0\arrow[r]&I\arrow[r]&H\otimes_A B\arrow[r]&M\otimes_A B\arrow[r]&Q\otimes_A B\arrow[r]&0
\end{tikzcd}}

\margin{53}
